% !TeX program = xelatex
\documentclass[11pt,a4paper]{article}
\usepackage{fontspec}
\usepackage{polyglossia}
\setdefaultlanguage{ukrainian}
\setmainfont{DejaVu Serif}
\setsansfont{DejaVu Sans}
\newfontfamily\cyrillicfont{DejaVu Serif}
\newfontfamily\cyrillicfontsf{DejaVu Sans}
\usepackage{amsmath,amssymb,booktabs,geometry,xcolor,titlesec}
\geometry{margin=23mm}
\definecolor{accent}{HTML}{1E685D}
\titleformat{\section}{\Large\sffamily\bfseries\color{accent}}{\thesection.}{0.5em}{}
\setlength{\parindent}{0pt}
\setlength{\parskip}{5pt}
\setlength{\emergencystretch}{2em}
\newcommand{\R}{\mathbb R}
\newcommand{\PP}{\mathcal P}
\begin{document}
{\sffamily\small ОБЧИСЛЮВАЛЬНА МАТЕМАТИКА ТА ПРОГРАМУВАННЯ \hfill ЛЕКЦІЯ 6}\\[5mm]
{\sffamily\huge Апроксимація табличних функцій}\\[2mm]
{\sffamily Математичний матеріал до лабораторної роботи №5}
\section{Табличні дані та проміжні значення}
Нехай $N\ge 2$, $x_0<x_1<\dots<x_{N-1}$, $y_i\in\R$.
Табличну функцію задає набір пар $(x_i,y_i)$, $i=0,\ldots,N-1$.
Інтерполянт $s$ відтворює значення в усіх вузлах:
\[
    s(x_i)=y_i,\quad i=0,\ldots,N-1.
\]
Значення $s(x_*)$ усередині $[x_0,x_{N-1}]$ є інтерполяційною оцінкою.
Застосування моделі за межами цього відрізка є екстраполяцією.
Без обмеження класу функцій інтерполянт не є єдиним.

\textbf{Лінійна модель на одному інтервалі.}
Для $x_i\le x\le x_{i+1}$, $h_i=x_{i+1}-x_i>0$,
\[
t=\frac{x-x_i}{h_i},\qquad
s(x)=y_i+\frac{y_{i+1}-y_i}{h_i}(x-x_i)
=(1-t)y_i+ty_{i+1}.
\]
Ваги $1-t$ і $t$ невід’ємні та мають суму $1$. Звідси
$\min(y_i,y_{i+1})\le s(x)\le\max(y_i,y_{i+1})$.
Якщо $y_i=f(x_i)$ і $f\in C^2[x_i,x_{i+1}]$, то
\[
|f(x)-s(x)|\le \frac{h_i^2}{8}\max_{\xi\in[x_i,x_{i+1}]}|f''(\xi)|.
\]
Межа стосується точних вузлових значень і містить припущення про гладкість $f$.

\section{Глобальна поліноміальна інтерполяція}
Позначимо через $\PP_d$ простір дійсних поліномів степеня не вище $d$.
Для $N$ попарно різних вузлів існує єдиний $p\in\PP_{N-1}$ з умовами $p(x_i)=y_i$.
Запишемо $p(x)=\sum_{j=0}^{N-1}a_jx^j$. Коефіцієнти визначає система
\[
 V a=y,\qquad V_{ij}=x_i^j,\quad 0\le i,j<N,\qquad
 \det V=\prod_{0\le i<j<N}(x_j-x_i)\ne0 .
\]
За зростанням степенів вектор $a=(a_0,\dots,a_{N-1})^\mathsf{T}$.
У записі від старшого степеня до вільного порядок коефіцієнтів обернений.
За $N=6$ маємо $p(x)=a_5x^5+a_4x^4+a_3x^3+a_2x^2+a_1x+a_0$.

\textbf{Еквівалентні формули.}
Форма Лагранжа:
\[
 L_i(x)=\prod_{\substack{0\le j<N\\j\ne i}}\frac{x-x_j}{x_i-x_j},
 \qquad L_i(x_j)=\delta_{ij},\qquad p(x)=\sum_{i=0}^{N-1}y_iL_i(x).
\]
При зміні лише $y_j$ на $\Delta y_j$ зміна полінома в довільній точці становить
$\Delta p(x)=\Delta y_j L_j(x)$.
У формі Ньютона з розділеними різницями:
\[
 f[x_i]=y_i,\qquad
 f[x_i,\ldots,x_{i+k}]
 =\frac{f[x_{i+1},\ldots,x_{i+k}]-f[x_i,\ldots,x_{i+k-1}]}{x_{i+k}-x_i},
\]
\[
 p(x)=f[x_0]+\sum_{k=1}^{N-1} f[x_0,\ldots,x_k]\prod_{j=0}^{k-1}(x-x_j).
\]
Для вузлів $(0,1),(1,2),(2,1)$:
\[
c=1,\quad a+b=1,\quad 4a+2b=0,\qquad p(x)=-x^2+2x+1,\quad p(1/2)=7/4.
\]
Контроль відтворення даних: $\max_i|p(x_i)-y_i|$.
У точній арифметиці ця величина нульова.
Вона не оцінює похибку між вузлами.

\section{Багатоінтервальна інтерполяція}
Для розбиття на інтервали з межами $\tau_0<\dots<\tau_m$ сплайн степеня $k$ є кусковим поліномом: на кожному інтервалі він належить $\PP_k$.
Для простих внутрішніх стиків $\tau_j$ і стандартного гладкого сплайна степеня $k\ge1$ накладають
\[
s_{j-1}^{(\ell)}(\tau_j)=s_{j}^{(\ell)}(\tau_j),
\qquad \ell=0,\ldots,k-1.
\]
Інтерполяційні умови $s(x_i)=y_i$ задають окремо.
Внутрішні стики сплайна та абсциси табличних даних можуть відрізнятися.

\begin{center}
\begin{tabular}{cll}
\toprule
$k$ & Локальна модель & Гладкість на простих стиках\\
\midrule
0 & Кускова стала & Допускає стрибки\\
1 & Лінійна & $C^0$\\
2 & Квадратична & $C^1$\\
3 & Кубічна & $C^2$\\
\bottomrule
\end{tabular}
\end{center}
Для утримання попереднього значення $s(x)=y_i$ за $x_i\le x<x_{i+1}$,
а $s(x_{N-1})=y_{N-1}$.
Для $k=1$ застосовується лінійна формула на кожній парі сусідніх вузлів.

\textbf{Квадратична модель зі стиками в табличних вузлах.}
На $[x_i,x_{i+1}]$ нехай $u=x-x_i$, $s_i=y_i+b_i u+c_i u^2$.
Інтерполяційна умова справа та неперервність нахилу дають
\[
c_i=\frac{y_{i+1}-y_i-b_i h_i}{h_i^2},\qquad b_{i+1}=b_i+2c_i h_i .
\]
Початковий нахил $b_0$ завершує цю конкретну постановку.
Інший вибір стиків чи додаткової умови може дати інший квадратичний інтерполянт.

\textbf{Кубічна модель зі стиками в табличних вузлах.}
На $[x_i,x_{i+1}]$ маємо $s_i=A_i+B_i u+C_i u^2+D_i u^3$.
Задають $s_i(x_i)=y_i$, $s_i(x_{i+1})=y_{i+1}$, неперервність $s'$ і $s''$ у внутрішніх вузлах та дві крайові умови.
Для природного сплайна $s''(x_0)=s''(x_{N-1})=0$.
Для умови \textit{not-a-knot} при $N\ge4$ перші два та останні два шматки попарно є обмеженнями одного кубічного полінома, що еквівалентно відповідним рівностям третіх похідних. При $N=4$ усі шматки належать одному поліному.
Кубічна інтерполяція загалом не зберігає межі $\min_i y_i$ і $\max_i y_i$.
Для вузлів $(0,0),(1,20),(2,80),(3,100)$ кубічний інтерполянт з умовою \textit{not-a-knot} має вигляд
\[
s(x)=-\tfrac{40}{3}x^3+60x^2-\tfrac{80}{3}x,\qquad s(1/4)=-25/8 .
\]

\section{Метод найменших квадратів}
Нехай задано $N$ вимірів $(x_i,y_i)$ та степінь $d$.
Для $p_a(x)=\sum_{j=0}^d a_j x^j$ визначимо залишки та критерій
\[
r_i=y_i-p_a(x_i),\qquad
S_d(a)=\sum_{i=0}^{N-1}r_i^2,\qquad
a^*\in\underset{a\in\R^{d+1}}{\operatorname{arg\,min}}\ S_d(a).
\]
За однакової ваги спостережень мінімізація $S_d$ еквівалентна мінімізації $S_d/N$ або $\sqrt{S_d/N}$.
Для матриці $X_{ij}=x_i^j$, $j=0,\dots,d$ маємо
\[
 S_d(a)=\|y-Xa\|_2^2,\qquad X^\mathsf{T}X a^*=X^\mathsf{T}y .
\]
Мінімізатор єдиний, якщо $\operatorname{rank}(X)=d+1$.
Для степеневого базису достатньо щонайменше $d+1$ різних абсцис.
Нормальні рівняння характеризують розв’язок; у чисельних обчисленнях придатні також QR-розклад або сингулярний розклад матриці $X$.

Для сталої $p(x)=\mu$ виконується тотожність
\[
  \sum_i(y_i-\mu)^2=\sum_i(y_i-\bar y)^2+N(\mu-\bar y)^2,\qquad
  \bar y=N^{-1}\sum_i y_i.
\]
Для $y=(99,100,101)$ отримуємо $S_0(\mu)=2+3(\mu-100)^2$.
За фіксованої таблиці $S_{d+1}(a^*_{d+1})\le S_d(a^*_d)$, бо $\PP_d\subset\PP_{d+1}$.
Ця нерівність стосується навчальної вибірки. Для оцінки прогнозу в нових точках потрібні окремі контрольні дані або додаткові припущення про функцію та похибки вимірів.
\end{document}
